\section{A half-moment of the critical branching random walk}

The branching random walk of the first section has Gaussian steps along a binary tree. A different critical tree, with steps of size $\pm1$ and Poisson offspring, has a half-moment small enough to see the typical tree rather than the rare tall one. The same half-moment is the model for the low moments of a random multiplicative function.

\section{The model: a critical branching random walk}

Think of the scales of primes as generations. Group the primes into blocks $(e^{e^{k-1}},e^{e^k}]$; each block contributes to $\log|F(\frac12+it)|$ a term of variance about $\frac12$, by Mertens's theorem. A point $t$ at resolution $e^{-k}$ is determined by its first $k$ blocks, and two points at distance about $e^{-j}$ share the first $j$ blocks and then go their separate ways: that is a tree of depth $n\approx\log\log x$, with independent increments along its edges. The following model keeps exactly this structure and nothing else.

\begin{definition}\label{def:brw}
In the \emph{critical branching random walk}, every individual, independently, has a random number $N^+\sim\mathrm{Poisson}(\frac1{2e})$ of children placed one step up and $N^-\sim\mathrm{Poisson}(\frac e2)$ of children placed one step down. The ancestor sits at $0$. For an individual $v$ of generation $|v|=n$ write $V(v)$ for its position, $v_k$ for its ancestor in generation $k$, and put
\[
W_n=\sum_{|v|=n}e^{V(v)}.
\]
\end{definition}

The numbers $\frac1{2e}$ and $\frac e2$ are chosen so that the two conditions of criticality hold: the expected sum of $e^{\text{step}}$ over the children of one individual is $\frac1{2e}\,e+\frac e2\,e^{-1}=1$, and the expected sum of $\text{step}\cdot e^{\text{step}}$ is $\frac12-\frac12=0$. The first makes $W_n$ a martingale with $\mathbb EW_n=1$; the second puts it at the critical point. (The mean number of children is $\frac1{2e}+\frac e2\approx1.54$, so the population grows exponentially when it survives.) The quantity $W_n$ plays the role of $V(x)/x$, and $n$ the role of $\log\log x$.

\begin{lemma}[many-to-one]\label{lem:manytoone}
For every $n\ge1$ and every function $G:\Z^n\to[0,\infty)$,
\[
\mathbb E\sum_{|v|=n}e^{V(v)}G\bigl(V(v_1),\dots,V(v_n)\bigr)=\mathbb E\,G(S_1,\dots,S_n),
\]
where $(S_k)$ is the simple random walk started at $0$.
\end{lemma}

\begin{proof}
Induction on $n$. Condition on the first generation. Given its individuals and their positions $V(c)=\pm1$, the subtrees rooted at them are independent copies of the whole process. The descendants of $c$ in generation $n$ are the individuals $w$ of generation $n-1$ of its subtree, with $V(v)=V(c)+V'(w)$, so $e^{V(v)}=e^{V(c)}e^{V'(w)}$. Let
\[
H(y)=\mathbb E\,G(y,y+S_1,\dots,y+S_{n-1})
\]
(for $n=1$, $H(y)=G(y)$). By the induction hypothesis applied to the subtree of $c$, the conditional expectation of the contribution of the descendants of $c$ is $e^{V(c)}H(V(c))$. It remains to compute $\mathbb E\sum_{|c|=1}e^{V(c)}H(V(c))$. A Poisson number of children with mean $m$ contributes $m$ times the value in expectation, so this is
\[
\tfrac1{2e}\,e\,H(1)+\tfrac e2\,e^{-1}H(-1)=\tfrac12H(1)+\tfrac12H(-1)=\mathbb E\,H(S_1),
\]
which is $\mathbb E\,G(S_1,\dots,S_n)$ by the Markov property of the walk.
\end{proof}

With $G=1$ the lemma gives $\mathbb EW_n=1$. Two more facts about the simple random walk are needed, both consequences of the reflection principle for the simple symmetric random walk on $\Z$.

\begin{lemma}[ballot bound]\label{lem:ballotbound}
For integers $L\ge0$ and $n\ge1$, $P(S_k\le L\text{ for all }k\le n)\le\dfrac{2L+2}{\sqrt n}$.
\end{lemma}

\begin{proof}
For the simple symmetric random walk and an integer $r\ge1$, $P(\max_{k\le n}S_k\ge r)=P(S_n\ge r)+P(S_n>r)$: every path that touches $r$ has a first time $\tau$ there, and reflecting the path up to $\tau$ is a bijection between $\{\max\ge r,\,S_n\ge r\}$ and $\{\max\ge r,\,S_n<r\}$, while $\{S_n>r\}\subset\{\max\ge r\}$. With $r=L+1$ and the symmetry of $S_n$,
\[
P\bigl(\max_{k\le n}S_k\le L\bigr)=1-P(S_n\ge L+1)-P(S_n\ge L+2)=P(-L-1\le S_n\le L),
\]
which is a sum of $2L+2$ point probabilities. Each is at most $a_n=\binom n{\lfloor n/2\rfloor}2^{-n}$. We show $a_n\le n^{-1/2}$. Let $b_m=\binom{2m}m4^{-m}$; then $b_{m+1}=b_m\frac{2m+1}{2m+2}$, and $b_m^2\le\frac1{3m+1}$ by induction, since $b_0=1$ and
\[
\Bigl(\frac{2m+1}{2m+2}\Bigr)^2\le\frac{3m+1}{3m+4}\iff12m^3+28m^2+19m+4\le12m^3+28m^2+20m+4.
\]
For $n=2m$, $a_n=b_m\le(3m+1)^{-1/2}\le n^{-1/2}$. For $n=2m+1$, $a_n=\frac12\binom{2m+1}m4^{-m}=b_m\frac{2m+1}{2m+2}\le b_m\le(3m+1)^{-1/2}\le n^{-1/2}$.
\end{proof}

\begin{lemma}[no high individuals]\label{lem:nohigh}
Let $M_n=\max\{V(v):|v|\le n\}$. For every integer $L\ge0$, $P(M_n\ge L+1)\le e^{-(L+1)}$.
\end{lemma}

\begin{proof}
If $M_n\ge L+1$, some individual $v$ with $|v|=k\le n$ is the \emph{first} on its line of descent to reach $L+1$: $V(v)=L+1$ (the walk moves by $\pm1$) and $V(v_j)\le L$ for $j<k$. Let $N$ be the number of such individuals. By Lemma~\ref{lem:manytoone} with $G=e^{-y_k}\mathbf 1\{y_k=L+1,\ y_j\le L\ (j<k)\}$,
\[
\mathbb EN=\sum_{k=1}^n e^{-(L+1)}P\bigl(S_k=L+1,\ S_j\le L\ (j<k)\bigr)\le e^{-(L+1)},
\]
since the events for different $k$ are disjoint. So $P(M_n\ge L+1)\le P(N\ge1)\le\mathbb EN$.
\end{proof}

\begin{theorem}[better than square root, in the tree]\label{thm:halfmoment}
For every $n\ge1$, $\mathbb E\bigl[W_n^{1/2}\bigr]\le\dfrac{C}{n^{1/4}}$, with $C=\sum_{L\ge0}\sqrt{2L+2}\,e^{-L/2}<5.4$. In particular $W_n\to0$ in probability, although $\mathbb EW_n=1$ for every $n$.
\end{theorem}

\begin{proof}
Since the ancestor sits at $0$, $M_n$ is an integer $\ge0$, and
\[
\mathbb E\bigl[W_n^{1/2}\bigr]=\sum_{L\ge0}\mathbb E\bigl[W_n^{1/2}\mathbf 1\{M_n=L\}\bigr].
\]
On the event $\{M_n=L\}$, every line of descent stays at or below $L$, so $W_n=W_n^{(L)}$, where $W_n^{(L)}=\sum_{|v|=n}e^{V(v)}\mathbf 1\{V(v_k)\le L\text{ for all }k\le n\}$. By the Cauchy--Schwarz inequality,
\[
\mathbb E\bigl[W_n^{1/2}\mathbf 1\{M_n=L\}\bigr]\le\bigl(\mathbb EW_n^{(L)}\bigr)^{1/2}P(M_n=L)^{1/2}.
\]
By Lemma~\ref{lem:manytoone} and Lemma~\ref{lem:ballotbound}, $\mathbb EW_n^{(L)}=P(S_k\le L\ \forall k\le n)\le\frac{2L+2}{\sqrt n}$. By Lemma~\ref{lem:nohigh}, $P(M_n=L)\le P(M_n\ge L)\le e^{-L}$ (trivially for $L=0$). Summing,
\[
\mathbb E\bigl[W_n^{1/2}\bigr]\le\sum_{L\ge0}\Bigl(\frac{2L+2}{\sqrt n}\Bigr)^{1/2}e^{-L/2}=\frac C{n^{1/4}}.
\]
The series converges, and its value is below $5.4$ (a direct computation gives $5.39\ldots$). Finally, by Markov's inequality, $P(W_n>\varepsilon)\le\varepsilon^{-1/2}\mathbb E[W_n^{1/2}]\to0$.
\end{proof}

Look at what the proof says. The mean $\mathbb EW_n=1$ is carried by rare trees in which some line of descent climbs high: there $W_n$ is huge. On the typical tree, all lines stay below a barrier at height $O(1)$, and the ballot theorem says that only a proportion $n^{-1/2}$ of the mass survives such a constraint. The half moment sees the typical tree, so it sees $n^{-1/4}$. Figure~\ref{fig:barrier} shows a sample.

\begin{figure}[ht]\centering\fig{ch06-barrier}
\caption{A critical branching random walk over $11$ generations. The highest individual ever reached height $M=2$ (dashed); the four lineages carrying most of $W_{11}=\sum e^{V(v)}$ are highlighted. Lineages that climb above the barrier are rare, and in the typical tree the mass that survives below it is a proportion $n^{-1/2}$ of the mean.}\label{fig:barrier}
\end{figure}

\begin{remark}[the lower bound, and the truth]\label{rem:brwlower}
The bound of Theorem~\ref{thm:halfmoment} is sharp: $\mathbb E[W_n^{1/2}]\ge cn^{-1/4}$ for some $c>0$. The matching lower bound applies $\mathbb EZ^{1/2}\ge(\mathbb EZ)^{3/2}/(\mathbb EZ^2)^{1/2}$ to a version of $W_n^{(0)}$ restricted to lines that end not too far below the barrier, and bounds $\mathbb EZ^2$ by a many-to-two lemma; it is longer and is not written here. Much more is true: $\sqrt n\,W_n$ converges in probability to a positive multiple of the limit of the \emph{derivative martingale} on the event of survival (Aïdékon and Shi, 2014, after Hu and Shi, 2009).
\end{remark}
